% !TeX root = SoulKnight_Report.tex

\documentclass[12pt,a4paper]{report}
\usepackage[utf8]{inputenc}
\usepackage[T5]{fontenc}
\usepackage[vietnamese]{babel}
\usepackage{listings}
\usepackage{xcolor}
\usepackage{geometry}

\geometry{
    left=2.5cm,
    right=2.5cm,
    top=2.5cm,
    bottom=2.5cm
}
\title{
    \Huge\textbf{SPIRIT KNIGHT}\\[0.5cm]
    \Large Technical Documentation
}

\author{Nguyễn Mạnh Cường \\
Vũ Việt Dũng \\
Trần Minh Tuấn Kiệt
}
\date{\today}
\begin{document}

    \maketitle
% tự động tạo mục lục
    \tableofcontents


    \chapter{Giới thiệu}


    \section{Tổng quan}

    Spirit Knight là một game nhập vai phát triển bằng Java và JavaFX.


    \section{Mục tiêu dự án}

    Mục tiêu của dự án là củng cố các kiến thức lập trình nâng cao, quản lí logic , xây dựng tư duy hướng đối tượng ,
    kĩ năng làm việc nhóm , sử dụng git


    \section{Sơ lược về cốt truyện và gameplay}
    Lấy cảm hứng từ tựa game Soul Knight - Chillyroom . Dự án được phát triển bởi nhóm sinh viên và học viên cao
    học ngành công nghệ thông tin . \\
    Bối cảnh bắt đầu từ một không gian quen thuộc với các Dev đó là một buổi đêm muộn đêm bàn máy.
    Trong cái đêm định mệnh ấy bạn nhận được một thông điệp lạ và lời thỉnh cầu đi giải cứu trái đất
    đang bị các thế lực đen tối dòm ngó. Bạn là người được chọn và hành trình của bạn sẽ bắt đầu
    sau cánh cổng portal


    \chapter{Kiến trúc hệ thống}


    \section{Tổng quan}
    Game được chia thành nhiều hệ thống như Player,Enemy,Weapon,Pet,Buff,Item,Database và UI


    \section{Hệ thống mã nguồn}
    Game được chia thành các package đảm bảo đưa các class liên quan về 1 package
    \begin{itemize}
        \item Package \texttt{animation}:gồm các class quản lí các hiệu ứng của game như:player/quái sinh/chết,
        render floating text , các hiệu ứng Particle pixel(các loai hạt), đổ bóng nhân vật
        \item  Package \texttt{buff}: gồm danh sách các buff của game,tính năng của từng loại trong các class
        riêng, Package nhỏ  \texttt{effect} phụ trách các hiệu ứng gây ra bởi các buff khi tấn công
        trực tiếp lên kẻ thù - các hiệu ứng của buff chủ yếu dùng render của javafx làm nên
        \item Package \texttt{cache}: phụ trách việc cache dữ liệu trong RAM, do game chạy bằng DB
        trực tuyến nên gây trễ khi chơi, việc cache sẽ phần nào giúp game chạy nhanh hơn và tăng trải nghiệm người chơi.
        \item Package \texttt{cinematic}: một package khá dặc biệt, nó góp phần tạo nên một phần intro story
        rất điện ảnh.
        \item Package \texttt{Database}: Linh hồn của dữ liệu game, lưu trữ hệ thống cơ sở dữ liệu của game
        sử dụng dịch vụ server của AVIEN dựa trên MySQL, tạo các truy vấn lên DB để lưu lại tên người chơi
        room họ đang chơi trước khi thoát,số điểm,số vàng họ kiếm được , các vật phẩm họ đã trang bị trong shop
        , các slot vũ khí họ sử dụng, tự động đồng bộ dữ liệu mỗi 15s
        \item Package \texttt{engine}: Đây là phần quan trọng nhất của game nơi chứa class \texttt{GameWorld.java}
        nơi gọi toàn bộ logic render,update player, enemy ,weapon, bullet, vật cản , pet , quản lí va chạm giữa các thành phần
        gọi các hiệu ứng của toàn bộ vòng lăp game.Đông thời chứa class \texttt{Camera.java} quản lí góc nhìn của game
        với trung tâm là player,class \texttt{InputHandler.java} phụ trách việc tiếp nhận dữ liệu đầu vào để điều
        khiển game -đặc biệt là có thể chơi cảm ứng với touchpad.
        \item Package \texttt{entity}: quản lí các đối tượng chính \textbf{Player} và \textbf{Enemy} kế thừa từ
        \textbf{Entity} của game.Quản lí các loại hero người chơi có thể lựa chọn,mua của Shop . Đồng thời
        cũng quản lí các loaij quái, logic quái tự điều khiển để đánh Player.
        \item Package \texttt{item}: quản lý các vật phẩm nhặt được trong màn chơi, trong đó có
        gold, gem, và buff
        \item Package \texttt{map}: phụ trách viếc load map từ file json, gán các chỉ số trong tiltmap với
        các assets, quản lí các \textbf{Room} và logic đóng mở cửa khi hoàn thành mission đánh quái
        ở từng phòng
        \item Package \texttt{pet}:Quản lí pet, các logic về pet- người bạn đồng hành trợ giúp Player
        đánh quái và các con pet để mua trong shop.
        \item Package \texttt{ui}:Quản lí toàn bộ giao diện game từ \textit{ intro,login,register} đến \textit{main menu
        shop,leaderboard,setting,HUD}.Đồng thời cũng đồng bộ dữ liệu với Database đặc trưng riêng cho từng
        người dùng.Trái tim là class \texttt{UIManager.java} gọi toàn bộ logic UI từ các file \textbf{controler} và
        các file giao diện \textbf{fxml} kết hợp với style trong các file \textbf{CSS} cung vơí DB
        quản lí các giao diện chuyên biệt với từng tài khoản khác nhau \textit{như chơi lần đầu, đợi load
        dữ liệu,...} .
        \item Package \texttt{utils}:Quản lí các tiện ích cơ bản như các hằng số , load âm thanh,DatabaseExcuter
        và vector2D
        \item Package \texttt{wapon}: quản lý hệ thống vũ khí , các loại vũ khí được laod lên shop và lopic chiến
        đấu của từng loại.
        \item Class \texttt{Main.java}
        là điểm khởi đầu (entry point) của ứng dụng
        Spirit Knight. Lớp này kế thừa \texttt{Application} của JavaFX
        và chịu trách nhiệm khởi tạo các thành phần chính cần thiết trước
        khi trò chơi bắt đầu hoạt động.\\
        Các trách nhiệm chính của lớp \texttt{Main} bao gồm:

        \begin{itemize}
            \item Khởi tạo cửa sổ JavaFX và vùng render của game.
            \item Khởi tạo hệ thống xử lý input.
            \item Khởi tạo \texttt{GameWorld}.
            \item Khởi tạo và liên kết \texttt{UIManager} với game.
            \item Khởi tạo và điều khiển \texttt{GameLoop}.
            \item Khởi tạo kết nối cơ sở dữ liệu trên một luồng riêng.
            \item Quản lý quá trình lưu dữ liệu và giải phóng tài nguyên
            khi người chơi thoát game.
        \end{itemize}
    \end{itemize}
    \section{Class Main.java}
    \subsection{Khởi tạo giao diện JavaFX}
    Khi ứng dụng được khởi chạy, JavaFX gọi phương thức
    \texttt{start(Stage stage)}. Đây là nơi lớp \texttt{Main}
    khởi tạo cửa sổ chính của trò chơi.
    Một đối tượng \texttt{Canvas} được tạo với kích thước mặc định
    được định nghĩa trong lớp \texttt{Constants}. Canvas đóng vai trò
    là bề mặt chính để render các thành phần của thế giới game.
%    cai nay de viet code de len
    \begin{lstlisting}[language=Java]
Canvas canvas = new Canvas(
    Constants.WINDOW_WIDTH,
    Constants.WINDOW_HEIGHT
);

GraphicsContext gc = canvas.getGraphicsContext2D();
    \end{lstlisting}

    Đối tượng \texttt{GraphicsContext} được lấy từ Canvas và được
    sử dụng bởi hệ thống render để vẽ các thành phần của trò chơi.

\subsection{Khởi tạo hệ thống Input}

Lớp \texttt{InputHandler} chịu trách nhiệm tiếp nhận các sự kiện
đầu vào từ người chơi như bàn phím và các thao tác điều khiển
được hỗ trợ bởi game.
InputHandler được liên kết với \texttt{Scene} để nhận các sự kiện
đầu vào và sau đó được truyền vào \texttt{GameWorld}.

\subsection{Khởi tạo GameWorld và UIManager}

Sau khi hệ thống input được thiết lập, \texttt{Main} khởi tạo
đối tượng \texttt{GameWorld}.

\begin{lstlisting}[language=Java]
GameWorld world = new GameWorld(inputHandler);
\end{lstlisting}

\texttt{GameWorld} đóng vai trò quản lý trạng thái và các thành
phần chính của thế giới game.

Tiếp theo, hệ thống giao diện được khởi tạo thông qua
\texttt{UIManager}.

\begin{lstlisting}[language=Java]
\end{lstlisting}

Việc liên kết \texttt{UIManager} với \texttt{GameWorld} cho phép
giao diện truy cập các trạng thái cần thiết của game để cập nhật
HUD và các thành phần giao diện khác.
    \subsection{Khởi tạo Game Loop}

    Lớp \texttt{Main} khởi tạo \texttt{GameLoop} để thực hiện vòng
    lặp chính của trò chơi. Mỗi vòng lặp nhận giá trị
    \texttt{delta}, biểu diễn khoảng thời gian giữa hai frame liên tiếp.

    Trong mỗi frame, hệ thống thực hiện ba nhiệm vụ chính:

    \begin{enumerate}
        \item Xử lý input toàn cục.
        \item Cập nhật trạng thái thế giới game.
        \item Render thế giới và cập nhật HUD.
    \end{enumerate}
    \subsection{Khởi tạo cơ sở dữ liệu bất đồng bộ}

    Việc kiểm tra và khởi tạo cơ sở dữ liệu được thực hiện trên một
    luồng riêng thay vì trực tiếp trên JavaFX Application Thread.
    Thiết kế này giúp quá trình kết nối cơ sở dữ liệu không chặn
    JavaFX Application Thread, từ đó hạn chế tình trạng giao diện
    bị đứng trong quá trình khởi động game.
    Luồng này được thiết lập là daemon thread để không ngăn JVM
    kết thúc khi ứng dụng được đóng.

    \chapter{Hệ thống GameWorld}
    \section{Tổng quan}
    Lớp \texttt{GameWorld} thuộc package
    \texttt{com.soulknight.engine} và đóng vai trò là bộ điều phối
    trung tâm của thế giới trò chơi Spirit Knight.
    Khác với lớp \texttt{Main} chịu trách nhiệm khởi tạo ứng dụng,
    \texttt{GameWorld} quản lý trạng thái của một phiên chơi và
    điều phối hoạt động giữa các hệ thống gameplay.
    Các thành phần được quản lý bởi \texttt{GameWorld} bao gồm:
    \begin{itemize}
        \item Player và Pet.
        \item Enemy và quá trình spawn Enemy.
        \item Map, Room và Obstacle.
        \item Bullet và hệ thống va chạm.
        \item Weapon và loadout của lượt chơi.
        \item Item và quá trình thu thập vật phẩm.
        \item Buff và các hiệu ứng chiến đấu.
        \item Particle và các hiệu ứng hình ảnh.
        \item Mission và Reward.
        \item Camera và quá trình render thế giới.
        \item Game State.
        \item Save/Load dữ liệu người chơi.
    \end{itemize}
    Có thể xem \texttt{GameWorld} như lớp trung gian kết nối các
    subsystem gameplay với vòng lặp chính của trò chơi.
    \section{Kiến trúc tổng thể}
    Về mặt kiến trúc, \texttt{GameWorld} không trực tiếp cài đặt
    toàn bộ hành vi của từng đối tượng. Thay vào đó, lớp giữ tham
    chiếu đến các manager và entity, sau đó điều phối chúng trong
    quá trình update và render.
    Mô hình tổng quát có thể biểu diễn như sau:

    \begin{verbatim}
                    GameWorld
                        |
       +----------------+----------------+
       |                |                |
       v                v                v
     Player         MapManager       LevelManager
       |                |                |
       v                v                v
     Weapon           Room         MissionManager
       |
       +----------+
       |          |
       v          v
     Bullet     Combat
                  |
        +---------+---------+
        |                   |
        v                   v
      Enemy               Effects

                    GameWorld
                        |
              +---------+---------+
              |                   |
              v                   v
          Save/Load             Render
              |                   |
              v                   v
           Database             Camera
    \end{verbatim}

    Thiết kế này cho phép các hệ thống chuyên biệt như
    \texttt{MapManager}, \texttt{LevelManager},
    \texttt{MissionManager} hay \texttt{CombatEffectManager}
    tập trung vào trách nhiệm riêng, trong khi
    \texttt{GameWorld} quyết định thời điểm chúng được cập nhật.
    \section{Các thành phần chính}
    \subsection{Các Manager}
    \texttt{GameWorld} sở hữu nhiều manager phục vụ cho các hệ thống
    khác nhau của game.
    \begin{table}[h]
        \centering
        \begin{tabular}{|l|p{8cm}|}
            \hline
            \textbf{Thành phần} & \textbf{Trách nhiệm} \\
            \hline
            \texttt{LevelManager}
            & Quản lý level và tiến trình của một lượt chơi. \\
            \hline
            \texttt{MapManager}
            & Quản lý bản đồ, tile, room, portal và dữ liệu không gian. \\
            \hline
            \texttt{MissionManager}
            & Quản lý nhiệm vụ của level hiện tại. \\
            \hline
            \texttt{EnemyFactory}
            & Khởi tạo các Enemy phù hợp với level. \\
            \hline
            \texttt{ParticleManager}
            & Quản lý các particle effect trong game. \\
            \hline
            \texttt{EnemyBurnManager}
            & Quản lý trạng thái Burn của Enemy. \\
            \hline
            \texttt{CombatEffectManager}
            & Quản lý các hiệu ứng chiến đấu đặc biệt. \\
            \hline
            \texttt{ItemMagnetSystem}
            & Điều khiển cơ chế hút Item về phía Player. \\
            \hline
            \texttt{FloatingTextManager}
            & Hiển thị damage, gold và các thông báo nổi trong game. \\
            \hline
        \end{tabular}
        \caption{Một số subsystem được GameWorld quản lý}
        \label{tab:gameworld-managers}
    \end{table}
    \section{Quản lý trạng thái trò chơi}
    Một trong những trách nhiệm quan trọng nhất của
    \texttt{GameWorld} là quản lý trạng thái hiện tại của trò chơi
    thông qua \texttt{GameState}.
    Các trạng thái được xử lý trong GameWorld bao gồm:
    \begin{itemize}
        \item \texttt{INTRO}
        \item \texttt{MAIN\_MENU}
        \item \texttt{PLAYING}
        \item \texttt{PAUSED}
        \item \texttt{LEVEL\_CLEAR}
        \item \texttt{REWARD\_PICK}
        \item \texttt{GAME\_OVER}
        \item \texttt{GAME\_VICTORY}
    \end{itemize}
    \subsection{Chuyển đổi trạng thái}
    Việc chuyển trạng thái được tập trung trong phương thức
    \texttt{changeState()}.
    Khi trạng thái thay đổi, trạng thái input cũ được xóa nhằm tránh
    việc một thao tác bàn phím hoặc chuột từ state trước tiếp tục
    ảnh hưởng tới state mới.
    Sau đó \texttt{GameStateListener} được thông báo để các hệ thống
    khác, đặc biệt là UI, có thể phản ứng với trạng thái mới.
    \section{Chu trình cập nhật GameWorld}
    \texttt{GameWorld} trong mỗi frame.
    Nếu trò chơi đang ở trạng thái \texttt{PAUSED}, phương thức
    kết thúc ngay lập tức và toàn bộ gameplay ngừng cập nhật.
    Đối với các trạng thái còn lại, GameWorld sử dụng
    \texttt{switch} để chuyển việc xử lý đến logic tương ứng.
    \begin{verbatim}
GameLoop
   |
   v
GameWorld.update()
   |
   +-- PAUSED ---------> Stop Update
   |
   +-- INTRO ----------> Input
   |
   +-- MAIN_MENU ------> Input
   |
   +-- PLAYING --------> updatePlaying()
   |
   +-- LEVEL_CLEAR ----> updateLevelClear()
   |
   +-- REWARD_PICK ----> updateRewardPick()
   |
   +-- GAME_OVER ------> Restart
   |
   +-- GAME_VICTORY ---> Restart
    \end{verbatim}
    \section{Cập nhật gameplay}
    Khi state hiện tại là \texttt{PLAYING}, GameWorld gọi phương
    thức \texttt{updatePlaying()}.
    Đây là một trong những phương thức điều phối quan trọng nhất
    của hệ thống gameplay.
    Quá trình update được thực hiện theo thứ tự tổng quát:
    \begin{enumerate}
        \item Cập nhật hiệu ứng spawn và death.
        \item Cập nhật Player.
        \item Cập nhật Pet.
        \item Xác định Room hiện tại.
        \item Cập nhật trạng thái Room.
        \item Cập nhật Enemy.
        \item Xử lý va chạm Player--Enemy.
        \item Cập nhật Bullet.
        \item Xử lý Obstacle bị phá hủy.
        \item Cập nhật các Combat Effect.
        \item Cập nhật trạng thái Burn.
        \item Cập nhật Item Magnet.
        \item Kiểm tra Player thu thập Item.
        \item Kiểm tra Reward.
        \item Thực hiện Auto Save.
        \item Cập nhật Camera.
        \item Kiểm tra trạng thái sống/chết của Player.
        \item Cập nhật Particle và Floating Text.
    \end{enumerate}
    Thứ tự update có ý nghĩa quan trọng vì nhiều subsystem phụ
    thuộc vào kết quả của subsystem trước đó.
    \section{Quản lý Player và Pet}
    \texttt{GameWorld} giữ tham chiếu đến Player hiên tại thông qua lời gọi đến class Player
    Trong trạng thái \texttt{PLAYING}, Player được cập nhật trước
    nhiều hệ thống chiến đấu khác.
    Việc truyền \texttt{this} vào \texttt{player.update()} cho phép
    Player tương tác với dữ liệu của thế giới như Enemy, Bullet,
    Map và các subsystem gameplay.
    \subsection{Pet System}
    Loại Pet được lấy từ \texttt{PetSelectionManager} và instance
    thực tế được tạo thông qua \texttt{PetFactory}.
    Nhờ đó, hệ thống lựa chọn Pet và đối tượng Pet tồn tại trong
    game được tách biệt với nhau.
    \section{Hệ thống chiến đấu}
    GameWorld đóng vai trò điều phối nhiều thành phần của Combat
    System, bao gồm:
    \begin{itemize}
        \item Bullet.
        \item Slash Effect.
        \item Explosion Effect.
        \item Ion Explosion.
        \item Sound Wave.
        \item Burn Effect.
        \item Damage.
        \item Combat Buff.
    \end{itemize}
    \subsection{Quản lý Bullet}
    Các projectile đang tồn tại được lưu trong:
    \begin{lstlisting}[language=Java]
private final List<Bullet> bullets = new ArrayList<>();
    \end{lstlisting}
    Trong mỗi frame, \texttt{updateBullets()} chịu trách nhiệm cập
    nhật vị trí và kiểm tra va chạm.
    Một Bullet có thể tương tác với:
    \begin{itemize}
        \item Wall.
        \item Obstacle.
        \item Enemy.
        \item Player.
    \end{itemize}
    Tùy loại vũ khí, Bullet còn có thể có các hành vi đặc biệt như
    phản xạ, xuyên mục tiêu, tạo explosion hoặc tạo sound wave.
    \subsection{Damage}
    Khi Bullet của Player va chạm Enemy, GameWorld tính lượng damage
    thực tế dựa trên HP trước và sau khi nhận sát thương.
    Mô hình xử lý có dạng:
    \begin{verbatim}
Bullet
   |
   v
Collision
   |
   v
Enemy.takeDamage()
   |
   +----> Floating Damage Text
   |
   +----> Particle Effect
   |
   +----> Buff Notification
    \end{verbatim}
    Việc sử dụng lượng damage thực tế giúp hệ thống Buff nhận đúng
    lượng damage đã gây ra thay vì chỉ sử dụng damage lý thuyết
    của projectile.
    \section{Hệ thống kỹ năng và hiệu ứng chiến đấu}
    GameWorld còn điều phối các kỹ năng đặc biệt của Player-được kich hoạt khi sử dụng các Buffs
    Ví dụ, hệ thống hiện tại bao gồm:
    \begin{itemize}
        \item Dragon Breath.
        \item Holy Nova.
        \item Chain Lightning.
        \item Death Explosion.
        \item Ion Explosion.
    \end{itemize}
    Đối với Dragon Breath, GameWorld xác định Enemy nằm trong vùng
    hình nón trước mặt Player. Damage được xử lý trực tiếp bởi logic
    gameplay, trong khi animation chỉ đóng vai trò biểu diễn trực
    quan.
    Cách tách gameplay logic và visual effect này giúp kết quả combat
    không phụ thuộc vào thời gian chạy animation.
    \section{Quản lý Item}
    GameWorld sử dụng \texttt{ItemMagnetSystem} để cập nhật chuyển
    động hút Item về phía Player.
    Sau đó \texttt{updateItemCollection()} kiểm tra va chạm giữa
    Player và Item.
    Các loại Item hiện được xử lý bao gồm:
    \begin{itemize}
        \item Gold.
        \item Gem.
        \item Buff.
        \item Energy Crystal.
    \end{itemize}
    Sau khi Item được thu thập, GameWorld cập nhật dữ liệu tương ứng,
    tạo visual feedback và loại Item khỏi danh sách active.
    \section{Hệ thống Room và Reward}
    GameWorld theo dõi Room mà Player đang đứng
    Trạng thái Room ở frame trước cũng được lưu lại để GameWorld có
    thể phát hiện sự thay đổi trạng thái của phòng.
    Sau khi một phòng chiến đấu được hoàn thành, hệ thống có thể mở
    \texttt{RewardPicker}.
    Reward hiện tại cho phép Player lựa chọn giữa:
    \begin{itemize}
        \item Gold.
        \item Weapon mới.
    \end{itemize}
    Mỗi phòng chỉ được nhận reward một lần trong cùng một lượt chơi.
    \section{Hệ thống vũ khí trong một lượt chơi}
    GameWorld quản lý hai slot vũ khí riêng của lượt chơi
    Khi bắt đầu một run, loadout được snapshot từ
    \texttt{WeaponSelectionManager}.
    Nếu slot không có dữ liệu hợp lệ, hệ thống sử dụng vũ khí mặc
    định:
    \begin{itemize}
        \item Slot 1: \texttt{BLASTER}.
        \item Slot 2: \texttt{OLD\_SWORD}.
    \end{itemize}
    Thiết kế này tạo ra sự phân biệt giữa:
    \begin{verbatim}
Shop Loadout
     |
     | Start Run
     v
Run Loadout
     |
     +---- Slot 1
     |
     +---- Slot 2
    \end{verbatim}
    Do đó, việc nhận Weapon Reward trong dungeon có thể thay đổi
    vũ khí của run hiện tại mà không nhất thiết thay đổi lựa chọn
    vũ khí được lưu trong Shop.
    \section{Khởi tạo một lượt chơi}
    Khi bắt đầu một lượt chơi mới, GameWorld thực hiện quá trình
    khởi tạo thông qua \texttt{startNewRun()} và
    \texttt{loadCurrentLevel()}.
    Quá trình load level bao gồm:
    \begin{enumerate}
        \item Khởi tạo MapManager.
        \item Load Obstacle từ Tile Map.
        \item Xác định spawn point.
        \item Khởi tạo hoặc di chuyển Player.
        \item Khôi phục dữ liệu Save nếu có.
        \item Khởi tạo Pet.
        \item Tạo Spawn Effect.
        \item Trang bị Weapon.
        \item Xóa Entity và Effect của level trước.
        \item Tạo Mission.
        \item Khởi tạo Item đặc thù của level.
        \item Spawn Boss nếu level hiện tại là Boss Level.
    \end{enumerate}
    \section{Render thế giới}
    Phần render gameplay chính được thực hiện bởi
    \texttt{renderWorld()}.
    \subsection{Y-Sorting}
    Spirit Knight sử dụng góc nhìn giả 2.5D. Vì vậy, thứ tự render
    không thể chỉ phụ thuộc vào loại Entity.
    GameWorld tạo một danh sách các đối tượng render và gán cho mỗi
    đối tượng một giá trị độ sâu dựa trên tọa độ Y.
    Những đối tượng có vị trí thấp hơn trên màn hình sẽ được render
    sau, tạo cảm giác chúng nằm phía trước đối tượng có tọa độ Y
    nhỏ hơn.
    \subsection{Render Pipeline}
    Pipeline render tổng quát có thể biểu diễn:
    \begin{verbatim}
Dynamic Background
        |
        v
      Floor
        |
        v
Obstacle Shadows
        |
        v
+-----------------------+
|      Y-SORTING        |
|                       |
| Walls                 |
| Player                |
| Pet                   |
| Enemy                 |
| Obstacles             |
| Doors                 |
+-----------------------+
        |
        v
      Items
        |
        v
     Bullets
        |
        v
 Combat Effects
        |
        v
    Particles
        |
        v
 Floating Text
    \end{verbatim}
    \section{Camera}
    GameWorld sở hữu một đối tượng \texttt{Camera} dùng để chuyển
    đổi tọa độ thế giới sang tọa độ màn hình.
    Trong quá trình gameplay, camera liên tục follow Player:
    Ngoài ra, GameWorld kiểm tra object có nằm trong vùng camera
    hay không trước khi render. Điều này giúp tránh thực hiện các
    draw call không cần thiết đối với object nằm ngoài viewport.
    \section{Hệ thống Save và Load}
    GameWorld tích hợp trực tiếp với hệ thống persistence thông qua:
    \begin{itemize}
        \item \texttt{PlayerSave}.
        \item \texttt{PlayerSaveDAO}.
        \item \texttt{PlayerSaveMapper}.
        \item \texttt{ShopDAO}.
        \item \texttt{UserSession}.
    \end{itemize}
    \subsection{Database Worker}
    Các thao tác database được chuyển sang một luồng
    \texttt{ExecutorService} riêng.
    \begin{lstlisting}[language=Java]
private final ExecutorService databaseExecutor =
    Executors.newSingleThreadExecutor(runnable -> {
        Thread thread = new Thread(
            runnable,
            "soul-knight-database-worker"
        );

        thread.setDaemon(true);
        return thread;
    });
    \end{lstlisting}
    Mục tiêu của thiết kế này là tránh thực hiện các thao tác I/O
    database trực tiếp trên game thread, từ đó giảm nguy cơ gây
    khựng hoặc đứng game.
    \subsection{Cloud Save}
    Phương thức \texttt{saveGameAsync()} chuyển trạng thái hiện tại
    của GameWorld thành \texttt{PlayerSave} thông qua
    \texttt{PlayerSaveMapper}.
    Sau đó dữ liệu được lưu bằng \texttt{PlayerSaveDAO} trên
    database worker.
    \subsection{Cloud Load}
    Quá trình load được thực hiện bất đồng bộ thông qua
    \texttt{loadGameAsync()}.
    Dữ liệu lấy từ database được lưu tạm trong:
    \begin{lstlisting}[language=Java]
private PlayerSave pendingPlayerSave;
    \end{lstlisting}
    Sau khi Player và level đã được khởi tạo, dữ liệu này mới được
    áp dụng trở lại GameWorld.
    Cách thực hiện này tránh việc restore dữ liệu vào Player trước
    khi Player thực sự tồn tại.
    \section{Auto Save}
    GameWorld định nghĩa khoảng thời gian Auto Save:
    \begin{lstlisting}[language=Java]
private static final double AUTO_SAVE_INTERVAL = 30.0;
    \end{lstlisting}
    Trong quá trình gameplay, bộ đếm thời gian được cập nhật theo
    \texttt{deltaSeconds}. Khi đạt đến khoảng thời gian quy định,
    GameWorld yêu cầu lưu trạng thái hiện tại.
    Việc save thực tế vẫn được chuyển sang database worker để không
    làm gián đoạn game loop.
    \section{Quản lý Gold và Gem}
    GameWorld phân biệt tài nguyên thu được trong run với tài nguyên
    đã được đồng bộ vào tài khoản.
    Các biến:
    lưu lượng tài nguyên đang chờ đồng bộ.
    Khi Player nhận Gold hoặc Gem, giá trị tương ứng được cộng vào
    pending bank.
    Sau đó \texttt{syncBankRewardsAsync()} chuyển lượng tài nguyên
    này sang database.
    Nếu quá trình database thất bại, lượng Gold và Gem chưa đồng bộ
    được trả lại vào pending bank để có thể thử lại ở lần sau.
    Cơ chế này giúp giảm nguy cơ mất reward khi xảy ra lỗi kết nối
    database.
    \section{Quản lý đồng thời}
    Do GameWorld vừa hoạt động trên game thread vừa giao tiếp với
    database worker, một số trạng thái cần được bảo vệ khỏi việc
    truy cập đồng thời.
    Các đối tượng \texttt{AtomicBoolean} được sử dụng để ngăn nhiều
    thao tác Save, Load hoặc Bank Sync chạy đồng thời.
    Ngoài ra, dữ liệu reward đang chờ đồng bộ được bảo vệ bởi một
    lock riêng.
    Thiết kế này đặc biệt quan trọng vì database worker hoạt động
    trên thread khác với gameplay.
    \section{Quá trình Shutdown}

    Khi ứng dụng chuẩn bị đóng, phương thức
    \texttt{shutdown()} được sử dụng để kết thúc các tài nguyên
    thuộc GameWorld.
    Quá trình này bao gồm:
    \begin{enumerate}
        \item Đồng bộ Gold và Gem còn tồn đọng.
        \item Yêu cầu database executor dừng hoạt động.
    \end{enumerate}
    Điều này phối hợp với shutdown handler của lớp \texttt{Main}
    để hạn chế mất dữ liệu khi người chơi đóng game.
    \section{Luồng hoạt động tổng quát}
    Có thể tóm tắt vòng đời gameplay được điều phối bởi
    \texttt{GameWorld} như sau:
    \begin{verbatim}
                 MAIN MENU
                     |
              Start / Continue
                     |
                     v
              Load Game World
                     |
                     v
                  PLAYING
                     |
          +----------+----------+
          |                     |
          v                     v
       Update                 Render
          |                     |
          v                     v
 Player / Enemy / Map      Y-Sorting
 Bullet / Collision        Effects
 Item / Buff / Pet         Particles
 Reward / Mission          Floating Text
          |
          v
      Room Cleared
          |
          v
      Reward Pick
          |
          v
       PLAYING
          |
          v
      Level Clear
          |
       +--+--+
       |     |
       v     v
 Next Level  Victory
       |
       v
    Auto Save
    \end{verbatim}
    \section{Kết luận}
    \texttt{GameWorld} là một trong những lớp trung tâm của kiến trúc
    Spirit Knight. Lớp kết nối Game Loop với hầu hết các subsystem
    gameplay và chịu trách nhiệm duy trì trạng thái của một lượt chơi.
    Thông qua GameWorld, quá trình cập nhật Player, Enemy, Bullet,
    Room, Item, Pet, Mission, Reward, Camera, Render và Save/Load
    được thực hiện theo một trình tự thống nhất trong mỗi frame.
    Vì vậy, có thể xem \texttt{GameWorld} là lớp điều phối gameplay
    chính của Spirit Knight.







\end{document}